%%%%%%%%%%%%%%%%%%%%%%%%%%%%%%%%%%%%%%%%%%%%%%%%%%%%%%%
% CHAPTER: METHODOLOGY
%%%%%%%%%%%%%%%%%%%%%%%%%%%%%%%%%%%%%%%%%%%%%%%%%%%%%%%
\chapter{Methodology}

%%%%%%%%%%%%%%%%%%%%%%%%%%%%%%%%%%%%%%%%%%%%%%%%%%%%%%%
% METHOD: PROBLEM FORMULATION
%%%%%%%%%%%%%%%%%%%%%%%%%%%%%%%%%%%%%%%%%%%%%%%%%%%%%%%
\section{Problem Formulation}

Consider a nonlinear system with the following dynamics:
\begin{equation}
  \label{eq:system_dynamics}
  \ddtt\mv{x}
  =
  \mv{f}(\mv{x})
  +
  \mm{B}(\mv{x})
  \mysat
  \left(
    \mv{u}
  \right)
  ,
\end{equation}
where $\mv{x} \in \mathbb{R}^{n}$ is the state vector, $\mv{u} \in \mathbb{R}^{m}$ is the control input, $\mv{f}(\mv{x})$ is an unknown system function, $\mm{B}(\mv{x})$ is a unknown control effective matrix, and $\mysat(\cdot)$ represents the saturation function.

The control effectiveness matrix $\mm{B}:=\mm{B}(\mv{x})$ is assumed to be ...



%%%%%%%%%%%%%%%%%%%%%%%%%%%%%%%%%%%%%%%%%%%%%%%%%%%%%%%
% METHOD: PROPOSED APPROACH
%%%%%%%%%%%%%%%%%%%%%%%%%%%%%%%%%%%%%%%%%%%%%%%%%%%%%%%
\section{CONAC Formulation}

The DNN is defined as a function $\NN(\NNin;\wth)$

%%%%%%%%%%%%%%%%%%%%%%%%%%%%%%%%%%%%%%%%%%%%%%%%%%%%%%%
% METHOD: HARDNET++ FORMULATION
%%%%%%%%%%%%%%%%%%%%%%%%%%%%%%%%%%%%%%%%%%%%%%%%%%%%%%%
\section{HardNet++ Application}

%%%%%%%%%%%%%%%%%%%%%%%%%%%%%%%%%%%%%%%%%%%%%%%%%%%%%%%
% METHOD: STABILITY ANALYSIS
%%%%%%%%%%%%%%%%%%%%%%%%%%%%%%%%%%%%%%%%%%%%%%%%%%%%%%%
\section{Stability Analysis}



